%% HAND-WRITTEN fragment -- NOT generated by maestro2tex.
%% The TRM has no bibliography and no \cite anywhere, so the sources abbreviated
%% in Table~\ref{tab:electrode-refs} are spelled out here in prose. Keep this
%% list in step with the header of ic/castalia/electrode_bc/veriloga/veriloga.va.

\paragraph{Sources.} The abbreviations of Table~\ref{tab:electrode-refs} are:
B\&F, A.~J. Bard and L.~R. Faulkner, \emph{Electrochemical Methods:
Fundamentals and Applications}, 2nd ed., Wiley, 2001; B\&S, D.~Britz and
J.~Strutwolf, \emph{Digital Simulation in Electrochemistry}, 4th ed., Springer,
2016; Randles 1947, J.~E.~B. Randles, \emph{Discuss. Faraday Soc.} 1:11, 1947;
Feldberg 1969, S.~W. Feldberg, \emph{Electroanal. Chem.} 3:199, 1969;
Nicholson 1965, R.~S. Nicholson, \emph{Anal. Chem.} 37:1351, 1965;
Matsuda \& Ayabe 1955, H.~Matsuda and Y.~Ayabe, \emph{Z. Elektrochem.} 59:494,
1955. The peak-current law is due independently to Randles
(\emph{Trans. Faraday Soc.} 44:327, 1948) and \v{S}ev\v{c}\'ik
(\emph{Coll. Czech. Chem. Commun.} 13:349, 1948); the numerical form quoted in
Table~\ref{tab:electrode-refs} is the \SI{25}{\celsius} one of B\&F
Eq.~6.2.19.
